\chapter*{Further Reading}
\addcontentsline{toc}{chapter}{\protect\numberline{}Further Reading}
\markboth{FURTHER READING}{FURTHER READING}

This is not a bibliography. It is a list of the books that stand behind the one you just read, grouped by what they will do for you. If you want the mathematics at the level of theorems, Li and Vitányi's \textit{An Introduction to Kolmogorov Complexity and Its Applications} is the standard reference and far friendlier than its title; the algorithmic information theory in this book is treated there with care. The books below are for other kinds of depth.

\section*{The mathematics, at the source}

\textbf{Ray Solomonoff, ``A Formal Theory of Inductive Inference''} (1964). The two-part paper where the universal prior first appears. Dense, but the idea that starts everything in Part I is right there on the page: prediction as a weighting of every program by its length.

\textbf{Ming Li and Paul Vitányi, \textit{An Introduction to Kolmogorov Complexity and Its Applications}} (1993, with later editions). The standard text, and the one place to find every theorem of algorithmic information this book leans on: Kraft, monotone machines and the universal semimeasure, the invariance theorem, dominance. Rigorous and, against the odds, readable.

\textbf{Gregory Chaitin, \textit{Meta Math!}} (2005). The randomness side of the same mathematics, told by one of its inventors, with an enthusiasm that is its own argument.

\section*{Optimal minds}

\textbf{Marcus Hutter, \textit{Universal Artificial Intelligence}} (2005). Solomonoff induction turned into an agent: AIXI, an uncomputable ideal that pairs a universal mixture with expected-reward planning. Its optimality claims are real but depend on the chosen universal machine and on which notion of optimality you mean, a point Jan Leike and Hutter later made sharply. The technical depth behind what Chapter 5 only gestures at. Demanding, and worth it.

\section*{The computable multiverse}

\textbf{Max Tegmark, \textit{Our Mathematical Universe}} (2014). The fullest popular case for taking the multiverse of mathematical structures as real. Tegmark is warmer and more hopeful about it than this book; read him for the vision, then decide how much of it you will pay for.

\textbf{David Deutsch, \textit{The Fabric of Reality}} (1997). Computation, multiverse, and knowledge braided together by someone who thinks the branches are where everything important happens. Broad, confident, generative.

\textbf{Jürgen Schmidhuber, ``A Computer Scientist's View of Life, the Universe, and Everything''} (1997). The short paper that put the dovetailing generator on the table: the Library, run as a program. His speed prior, which charges for computation time as well as length, came a few years later.

\section*{The self, dissolved}

\textbf{Derek Parfit, \textit{Reasons and Persons}} (1984). A foundational treatment of personal identity and what matters when identity is not a further fact. Questions about copies and continuation arise here by another road. Chapter 14 is deeply indebted to it.

\textbf{Thomas Metzinger, \textit{The Ego Tunnel}} (2009). The self as a model the brain runs and cannot see through. Metzinger goes further than this book does. If Chapter 14 left you curious, he is the next step; if it left you unsettled, there is no hurry.

\section*{Living inside the structure}

\textbf{Albert Camus, \textit{The Myth of Sisyphus}} (1942). The source of the book's closing move. The indifference is given; the rock is given; the only question is what to do with the climb. \textit{Il faut imaginer Sisyphe heureux.}

\textbf{Simone Weil, \textit{Gravity and Grace}} (1947). Attention as a discipline and a gift; her line that attention is ``the rarest and purest form of generosity'' comes from a 1942 letter to Jo\"e Bousquet. Weil understood the indifference of the structure and chose to attend to the people inside it.

\section*{The Library in literature}

\textbf{Jorge Luis Borges, ``The Library of Babel'' and ``The Garden of Forking Paths''} (1941). The totality and the branching, as literature, decades before the mathematics named either. The first is this book's central image; the second belongs to \textit{Multitudes}. Each is short enough for one sitting.

\section*{The rest of the argument}

\textbf{Alex Towell, \textit{Worldlines: The Indifference of Geometry}} (2026) and \textbf{\textit{Multitudes: The Indifference of Measure}} (2026). The first two volumes, by the other two roads. Relativity and quantum mechanics, each read through one interpretation (the block universe, and Everett's), arrive near where this book arrives: an indifferent structure, and the warmth we choose to give each other inside it. Either can be read alone; the three together are one argument.
